% Options for packages loaded elsewhere
\PassOptionsToPackage{unicode}{hyperref}
\PassOptionsToPackage{hyphens}{url}
\PassOptionsToPackage{dvipsnames,svgnames,x11names}{xcolor}
%
\documentclass[
  ignorenonframetext,
  aspectratio=169,
]{beamer}
\usepackage{pgfpages}
\setbeamertemplate{caption}[numbered]
\setbeamertemplate{caption label separator}{: }
\setbeamercolor{caption name}{fg=normal text.fg}
\beamertemplatenavigationsymbolsempty
% Prevent slide breaks in the middle of a paragraph
\widowpenalties 1 10000
\raggedbottom
\setbeamertemplate{part page}{
  \centering
  \begin{beamercolorbox}[sep=16pt,center]{part title}
    \usebeamerfont{part title}\insertpart\par
  \end{beamercolorbox}
}
\setbeamertemplate{section page}{
  \centering
  \begin{beamercolorbox}[sep=12pt,center]{part title}
    \usebeamerfont{section title}\insertsection\par
  \end{beamercolorbox}
}
\setbeamertemplate{subsection page}{
  \centering
  \begin{beamercolorbox}[sep=8pt,center]{part title}
    \usebeamerfont{subsection title}\insertsubsection\par
  \end{beamercolorbox}
}
\AtBeginPart{
  \frame{\partpage}
}
\AtBeginSection{
  \ifbibliography
  \else
    \frame{\sectionpage}
  \fi
}
\AtBeginSubsection{
  \frame{\subsectionpage}
}
\usepackage{amsmath,amssymb}
\usepackage{iftex}
\ifPDFTeX
  \usepackage[T1]{fontenc}
  \usepackage[utf8]{inputenc}
  \usepackage{textcomp} % provide euro and other symbols
\else % if luatex or xetex
  \usepackage{unicode-math} % this also loads fontspec
  \defaultfontfeatures{Scale=MatchLowercase}
  \defaultfontfeatures[\rmfamily]{Ligatures=TeX,Scale=1}
\fi
\usepackage{lmodern}
\usetheme[]{metropolis}
\usefonttheme{professionalfonts}
\ifPDFTeX\else
  % xetex/luatex font selection
  \setmonofont[Scale=0.78]{DejaVu Sans Mono}
\fi
% Use upquote if available, for straight quotes in verbatim environments
\IfFileExists{upquote.sty}{\usepackage{upquote}}{}
\IfFileExists{microtype.sty}{% use microtype if available
  \usepackage[]{microtype}
  \UseMicrotypeSet[protrusion]{basicmath} % disable protrusion for tt fonts
}{}
\makeatletter
\@ifundefined{KOMAClassName}{% if non-KOMA class
  \IfFileExists{parskip.sty}{%
    \usepackage{parskip}
  }{% else
    \setlength{\parindent}{0pt}
    \setlength{\parskip}{6pt plus 2pt minus 1pt}}
}{% if KOMA class
  \KOMAoptions{parskip=half}}
\makeatother
\usepackage{xcolor}
\newif\ifbibliography
\usepackage{color}
\usepackage{fancyvrb}
\newcommand{\VerbBar}{|}
\newcommand{\VERB}{\Verb[commandchars=\\\{\}]}
\DefineVerbatimEnvironment{Highlighting}{Verbatim}{commandchars=\\\{\}}
% Add ',fontsize=\small' for more characters per line
\newenvironment{Shaded}{}{}
\newcommand{\AlertTok}[1]{\textcolor[rgb]{1.00,0.00,0.00}{\textbf{#1}}}
\newcommand{\AnnotationTok}[1]{\textcolor[rgb]{0.38,0.63,0.69}{\textbf{\textit{#1}}}}
\newcommand{\AttributeTok}[1]{\textcolor[rgb]{0.49,0.56,0.16}{#1}}
\newcommand{\BaseNTok}[1]{\textcolor[rgb]{0.25,0.63,0.44}{#1}}
\newcommand{\BuiltInTok}[1]{\textcolor[rgb]{0.00,0.50,0.00}{#1}}
\newcommand{\CharTok}[1]{\textcolor[rgb]{0.25,0.44,0.63}{#1}}
\newcommand{\CommentTok}[1]{\textcolor[rgb]{0.38,0.63,0.69}{\textit{#1}}}
\newcommand{\CommentVarTok}[1]{\textcolor[rgb]{0.38,0.63,0.69}{\textbf{\textit{#1}}}}
\newcommand{\ConstantTok}[1]{\textcolor[rgb]{0.53,0.00,0.00}{#1}}
\newcommand{\ControlFlowTok}[1]{\textcolor[rgb]{0.00,0.44,0.13}{\textbf{#1}}}
\newcommand{\DataTypeTok}[1]{\textcolor[rgb]{0.56,0.13,0.00}{#1}}
\newcommand{\DecValTok}[1]{\textcolor[rgb]{0.25,0.63,0.44}{#1}}
\newcommand{\DocumentationTok}[1]{\textcolor[rgb]{0.73,0.13,0.13}{\textit{#1}}}
\newcommand{\ErrorTok}[1]{\textcolor[rgb]{1.00,0.00,0.00}{\textbf{#1}}}
\newcommand{\ExtensionTok}[1]{#1}
\newcommand{\FloatTok}[1]{\textcolor[rgb]{0.25,0.63,0.44}{#1}}
\newcommand{\FunctionTok}[1]{\textcolor[rgb]{0.02,0.16,0.49}{#1}}
\newcommand{\ImportTok}[1]{\textcolor[rgb]{0.00,0.50,0.00}{\textbf{#1}}}
\newcommand{\InformationTok}[1]{\textcolor[rgb]{0.38,0.63,0.69}{\textbf{\textit{#1}}}}
\newcommand{\KeywordTok}[1]{\textcolor[rgb]{0.00,0.44,0.13}{\textbf{#1}}}
\newcommand{\NormalTok}[1]{#1}
\newcommand{\OperatorTok}[1]{\textcolor[rgb]{0.40,0.40,0.40}{#1}}
\newcommand{\OtherTok}[1]{\textcolor[rgb]{0.00,0.44,0.13}{#1}}
\newcommand{\PreprocessorTok}[1]{\textcolor[rgb]{0.74,0.48,0.00}{#1}}
\newcommand{\RegionMarkerTok}[1]{#1}
\newcommand{\SpecialCharTok}[1]{\textcolor[rgb]{0.25,0.44,0.63}{#1}}
\newcommand{\SpecialStringTok}[1]{\textcolor[rgb]{0.73,0.40,0.53}{#1}}
\newcommand{\StringTok}[1]{\textcolor[rgb]{0.25,0.44,0.63}{#1}}
\newcommand{\VariableTok}[1]{\textcolor[rgb]{0.10,0.09,0.49}{#1}}
\newcommand{\VerbatimStringTok}[1]{\textcolor[rgb]{0.25,0.44,0.63}{#1}}
\newcommand{\WarningTok}[1]{\textcolor[rgb]{0.38,0.63,0.69}{\textbf{\textit{#1}}}}
\usepackage{longtable,booktabs,array}
\usepackage{calc} % for calculating minipage widths
\usepackage{caption}
% Make caption package work with longtable
\makeatletter
\def\fnum@table{\tablename~\thetable}
\makeatother
\setlength{\emergencystretch}{3em} % prevent overfull lines
\providecommand{\tightlist}{%
  \setlength{\itemsep}{0pt}\setlength{\parskip}{0pt}}
\setcounter{secnumdepth}{-\maxdimen} % remove section numbering
\metroset{block=fill}
\setbeamertemplate{navigation symbols}{}
\usepackage{booktabs}
\usepackage{amsmath,amssymb}
\ifLuaTeX
  \usepackage{selnolig}  % disable illegal ligatures
\fi
\usepackage{bookmark}
\IfFileExists{xurl.sty}{\usepackage{xurl}}{} % add URL line breaks if available
\urlstyle{same}
\hypersetup{
  pdftitle={QPUBench},
  pdfauthor={MQS · mark@mqs.dk},
  colorlinks=true,
  linkcolor={Maroon},
  filecolor={Maroon},
  citecolor={Blue},
  urlcolor={Blue},
  pdfcreator={LaTeX via pandoc}}

\title{QPUBench}
\subtitle{Open Source Quantum Computing Benchmark Framework}
\author{MQS · mark@mqs.dk}
\date{2026}
\institute{Molecular Quantum Solutions}

\begin{document}
\frame{\titlepage}

\section{Framework Overview}\label{framework-overview}

\begin{frame}[fragile]{What is QPUBench?}
\phantomsection\label{what-is-qpubench}
\textbf{Modality-agnostic quantum benchmark framework}

\begin{itemize}
\tightlist
\item
  Separates \emph{what} you benchmark (schemas) from \emph{how}
  execution happens (adapters) and \emph{how results are stored}
  (stores)
\item
  Typed \textbf{Pydantic v2} schema layer --- zero quantum SDK
  dependencies in core
\item
  Gate-based circuits, MBQC on FPGA, molecular VQE, evolutionary circuit
  search, and error-mitigated runs all share one
  \texttt{BenchmarkRecord}
\item
  Schema \textbf{v1.12.0} --- 21 modules, language-agnostic JSON output
\end{itemize}

\vspace{0.5em}
\begin{block}{Core invariant}
\texttt{qpubench} itself never imports from any quantum library. Only adapters do.
\end{block}
\end{frame}

\begin{frame}[fragile]{Architecture}
\phantomsection\label{architecture}
\begin{columns}[T]
\begin{column}{0.5\textwidth}
\textbf{Core layer} \emph{(no SDK dependencies)}

\begin{itemize}
\tightlist
\item
  \texttt{schemas/} --- 21 Pydantic v2 modules
\item
  \texttt{runner.py} --- \texttt{BenchmarkRunner}
\item
  \texttt{store.py} --- \texttt{NDJSONStore}, \texttt{ParquetStore}
\end{itemize}
\end{column}

\begin{column}{0.5\textwidth}
\textbf{Integration layer} \emph{(your project)}

\begin{itemize}
\tightlist
\item
  \texttt{backends/} --- adapter protocols + stubs
\item
  \texttt{integrations/} --- copy into your project
\item
  \texttt{examples/} --- runnable demos
\end{itemize}
\end{column}
\end{columns}

\vspace{0.5em}

\begin{verbatim}
CircuitSpec / Problem  ──▶  Adapter  ──▶  BenchmarkRecord  ──▶  Store
\end{verbatim}
\end{frame}

\begin{frame}[fragile]{Key Design Principles}
\phantomsection\label{key-design-principles}
\begin{longtable}[]{@{}
  >{\raggedright\arraybackslash}p{(\columnwidth - 2\tabcolsep) * \real{0.4231}}
  >{\raggedright\arraybackslash}p{(\columnwidth - 2\tabcolsep) * \real{0.5769}}@{}}
\toprule\noalign{}
\begin{minipage}[b]{\linewidth}\raggedright
Principle
\end{minipage} & \begin{minipage}[b]{\linewidth}\raggedright
Implementation
\end{minipage} \\
\midrule\noalign{}
\endhead
\textbf{Schema-first} & Pydantic v2 is the stable contract --- all
records are valid JSON \\
\textbf{No SDK lock-in} & Core never imports any quantum library \\
\textbf{Three adapter protocols} & \texttt{BackendAdapter} ·
\texttt{AlgorithmAdapter} · \texttt{ErrorMitigationAdapter} \\
\textbf{Append-only stores} & NDJSON and Parquet for reproducible
sweeps \\
\textbf{Language-agnostic} & Any language can parse the NDJSON output \\
\bottomrule\noalign{}
\end{longtable}
\end{frame}

\begin{frame}[fragile]{Installation}
\phantomsection\label{installation}
\begin{Shaded}
\begin{Highlighting}[]
\CommentTok{\# Minimal — schemas + stubs, no quantum SDK required}
\ExtensionTok{pip}\NormalTok{ install .}

\CommentTok{\# With optional backends}
\ExtensionTok{pip}\NormalTok{ install }\StringTok{".[qiskit]"}     \CommentTok{\# Qiskit Aer + IBM Quantum Runtime V2}
\ExtensionTok{pip}\NormalTok{ install }\StringTok{".[storage]"}    \CommentTok{\# Parquet store (pyarrow + pandas)}
\ExtensionTok{pip}\NormalTok{ install }\StringTok{".[all]"}        \CommentTok{\# everything on PyPI}
\end{Highlighting}
\end{Shaded}

\textbf{Package managers}: \texttt{pip} · \texttt{uv} ·
\texttt{Poetry\ 2} · \texttt{conda}

\textbf{Credentials}: via \texttt{.env} file and
\texttt{BackendSpec.auth} (never hardcoded)
\end{frame}

\section{Schema Layer}\label{schema-layer}

\begin{frame}[fragile]{Core Schema Modules}
\phantomsection\label{core-schema-modules}
\begin{longtable}[]{@{}
  >{\raggedright\arraybackslash}p{(\columnwidth - 4\tabcolsep) * \real{0.2759}}
  >{\raggedright\arraybackslash}p{(\columnwidth - 4\tabcolsep) * \real{0.3448}}
  >{\raggedright\arraybackslash}p{(\columnwidth - 4\tabcolsep) * \real{0.3793}}@{}}
\toprule\noalign{}
\begin{minipage}[b]{\linewidth}\raggedright
Module
\end{minipage} & \begin{minipage}[b]{\linewidth}\raggedright
Modality
\end{minipage} & \begin{minipage}[b]{\linewidth}\raggedright
Key types
\end{minipage} \\
\midrule\noalign{}
\endhead
\texttt{primitives} & all & \texttt{QPUModality},
\texttt{CircuitFormat}, \texttt{PauliLabel}, \texttt{ComplexNumber} \\
\texttt{circuit} & all & \texttt{CircuitSpec} ---
\texttt{from\_openqasm3()}, \texttt{bind()},
\texttt{is\_parametric()} \\
\texttt{observable} & all & \texttt{SparsePauliObservable},
\texttt{PauliTerm} \\
\texttt{backend} & all & \texttt{BackendSpec} --- 27 factory
constructors \\
\texttt{execution} & all & \texttt{ExecutionOptions},
\texttt{AlgorithmSpec}, \texttt{ZNEConfig} \\
\texttt{error\_mitigation} & all & \texttt{ErrorMitigationStrategy},
provider configs, \texttt{QuantumAdvantageRecord} \\
\texttt{result} & all & \texttt{QuantumResult},
\texttt{ExpectationResult}, \texttt{ShotResult} \\
\texttt{record} & all & \texttt{BenchmarkRecord}, \texttt{VQAConfig} \\
\bottomrule\noalign{}
\end{longtable}
\end{frame}

\begin{frame}[fragile]{Domain-Specific Schema Modules}
\phantomsection\label{domain-specific-schema-modules}
\begin{longtable}[]{@{}
  >{\raggedright\arraybackslash}p{(\columnwidth - 2\tabcolsep) * \real{0.2759}}
  >{\raggedright\arraybackslash}p{(\columnwidth - 2\tabcolsep) * \real{0.7241}}@{}}
\toprule\noalign{}
\begin{minipage}[b]{\linewidth}\raggedright
Module
\end{minipage} & \begin{minipage}[b]{\linewidth}\raggedright
Modality / Framework
\end{minipage} \\
\midrule\noalign{}
\endhead
\texttt{mbqc} & MBQC-FPGA --- bit-exact 16-bit program word \\
\texttt{photonic} & Linear-optics chips, FBQC, HOM, photonic
VQE/analog \\
\texttt{gbs} & Gaussian Boson Sampling --- hafnian, vibronic,
TDM/Borealis \\
\texttt{qdk\_chemistry} & SCF → active space → QPE → resource
estimation \\
\texttt{qse} & Krylov Quantum Diagonalization (KQD / SQD) \\
\texttt{qesem} & QESEM (Qedma) --- noise scaling, QET, device
characterisation \\
\texttt{qcschema} & QCSchema / QCElemental / PennyLane
interoperability \\
\texttt{neutral\_atom} & Bloqade / Aquila AHS, Rydberg atom
arrangement \\
\texttt{slowquant} & SlowQuant UCC/VQE --- ansatz, SCF, linear
response \\
\texttt{error\_mitigation} & Q-CTRL · Mitiq · Haiqu · ParityQC · QMatter
· Quantum Motion · IBM V2 · Advantage Tracker \\
\bottomrule\noalign{}
\end{longtable}
\end{frame}

\begin{frame}[fragile]{QPU Modalities}
\phantomsection\label{qpu-modalities}
QPUBench natively covers \textbf{9 quantum computing paradigms}:

\begin{columns}[T]
\begin{column}{0.5\textwidth}
\begin{itemize}
\tightlist
\item
  \texttt{GATE\_BASED} --- qubit circuits (QASM2/3)
\item
  \texttt{MBQC} --- measurement-based QC on FPGA
\item
  \texttt{PHOTONIC\_LINEAR\_OPTICS} --- Fock-state chips
\item
  \texttt{FUSION\_BASED} --- FBQC resource states + fusion
\item
  \texttt{QPE} --- Quantum Phase Estimation / IQPE
\end{itemize}
\end{column}

\begin{column}{0.5\textwidth}
\begin{itemize}
\tightlist
\item
  \texttt{GBS} --- Gaussian Boson Sampling
\item
  \texttt{KQD} --- Krylov Quantum Diagonalization
\item
  \texttt{NEUTRAL\_ATOM} --- AHS on Rydberg atoms
\item
  \texttt{ANNEALING} --- quantum annealing
\end{itemize}
\end{column}
\end{columns}

\vspace{0.3em}

All modalities share the same \texttt{BenchmarkRecord} schema.
\end{frame}

\begin{frame}[fragile]{CircuitSpec --- One Type for All Circuits}
\phantomsection\label{circuitspec-one-type-for-all-circuits}
\begin{Shaded}
\begin{Highlighting}[]
\CommentTok{\# Gate{-}based (QASM2)}
\NormalTok{circuit }\OperatorTok{=}\NormalTok{ CircuitSpec(}
\NormalTok{    num\_qubits}\OperatorTok{=}\DecValTok{2}\NormalTok{,}
\NormalTok{    serialized}\OperatorTok{=}\StringTok{\textquotesingle{}OPENQASM 2.0;}\CharTok{\textbackslash{}n}\StringTok{include "qelib1.inc";}\CharTok{\textbackslash{}n}\StringTok{...\textquotesingle{}}\NormalTok{,}
\NormalTok{    observables}\OperatorTok{=}\NormalTok{[SparsePauliObservable(num\_qubits}\OperatorTok{=}\DecValTok{2}\NormalTok{, terms}\OperatorTok{=}\NormalTok{[}
\NormalTok{        PauliTerm(qubit\_indices}\OperatorTok{=}\NormalTok{(}\DecValTok{0}\NormalTok{, }\DecValTok{1}\NormalTok{),}
\NormalTok{                  pauli\_ops}\OperatorTok{=}\NormalTok{(PauliLabel.Z, PauliLabel.Z),}
\NormalTok{                  coefficient}\OperatorTok{=}\NormalTok{ComplexNumber(re}\OperatorTok{=}\FloatTok{1.0}\NormalTok{))}
\NormalTok{    ])],}
\NormalTok{)}

\CommentTok{\# OpenQASM 3.0 helper}
\NormalTok{circuit }\OperatorTok{=}\NormalTok{ CircuitSpec.from\_openqasm3(source, num\_qubits}\OperatorTok{=}\DecValTok{2}\NormalTok{)}

\CommentTok{\# Parametric VQE — bind() returns a new copy}
\NormalTok{ansatz }\OperatorTok{=}\NormalTok{ CircuitSpec(num\_qubits}\OperatorTok{=}\DecValTok{2}\NormalTok{, parameters}\OperatorTok{=}\NormalTok{[}\StringTok{"theta"}\NormalTok{], ...)}
\NormalTok{bound  }\OperatorTok{=}\NormalTok{ ansatz.bind(\{}\StringTok{"theta"}\NormalTok{: }\FloatTok{1.2566}\NormalTok{\})}

\CommentTok{\# Algorithm{-}driven (chemistry)}
\NormalTok{mol }\OperatorTok{=}\NormalTok{ CircuitSpec(}\BuiltInTok{format}\OperatorTok{=}\NormalTok{CircuitFormat.MOLECULE\_JSON,}
\NormalTok{                  serialized}\OperatorTok{=}\StringTok{"/path/to/He{-}ccpvdz.json"}\NormalTok{, num\_qubits}\OperatorTok{=}\DecValTok{0}\NormalTok{)}
\end{Highlighting}
\end{Shaded}
\end{frame}

\section{Adapters}\label{adapters}

\begin{frame}[fragile]{Three Adapter Protocols}
\phantomsection\label{three-adapter-protocols}
\begin{columns}[T]
\begin{column}{0.33\textwidth}
\textbf{\texttt{BackendAdapter}} \emph{Circuit in, result out}

\begin{verbatim}
CircuitSpec
    ↓ run()
QuantumResult
\end{verbatim}

Aer · Qrack · IBM IQM · MBQC-FPGA
\end{column}

\begin{column}{0.33\textwidth}
\textbf{\texttt{AlgorithmAdapter}} \emph{Problem spec in}

\begin{verbatim}
CircuitSpec
 (problem)
    ↓ run_algorithm()
(Result, VQAConfig)
\end{verbatim}

QForte · OpenFermion
\end{column}

\begin{column}{0.33\textwidth}
\textbf{\texttt{ErrorMitigationAdapter}} \emph{Wraps a BackendAdapter}

\begin{verbatim}
inner adapter
    ↓ run()
mitigate/suppress
    ↓
QuantumResult
\end{verbatim}

Fire Opal · Mitiq Haiqu · ParityQC
\end{column}
\end{columns}

\vspace{0.3em}

\texttt{BenchmarkRunner} dispatches automatically ---
\texttt{ErrorMitigationAdapter} registers like any
\texttt{BackendAdapter}.
\end{frame}

\begin{frame}[fragile]{BackendAdapter --- Estimator vs Sampler}
\phantomsection\label{backendadapter-estimator-vs-sampler}
\begin{Shaded}
\begin{Highlighting}[]
\KeywordTok{class}\NormalTok{ MyBackendAdapter:}
    \AttributeTok{@property}
    \KeywordTok{def}\NormalTok{ spec(}\VariableTok{self}\NormalTok{) }\OperatorTok{{-}\textgreater{}}\NormalTok{ BackendSpec: ...          }\CommentTok{\# hardware description}

    \KeywordTok{def}\NormalTok{ validate(}\VariableTok{self}\NormalTok{, circuit: CircuitSpec) }\OperatorTok{{-}\textgreater{}} \BuiltInTok{list}\NormalTok{[}\BuiltInTok{str}\NormalTok{]: ...}

    \KeywordTok{def}\NormalTok{ run(}\VariableTok{self}\NormalTok{, circuit: CircuitSpec,}
\NormalTok{            options: ExecutionOptions) }\OperatorTok{{-}\textgreater{}}\NormalTok{ QuantumResult:}
        \ControlFlowTok{if}\NormalTok{ circuit.observables:}
            \CommentTok{\# Estimator path → expectation\_values (VQE, QAOA)}
\NormalTok{            evs }\OperatorTok{=}\NormalTok{ [measure\_expectation(circuit, obs) }\ControlFlowTok{for}\NormalTok{ obs }\KeywordTok{in}\NormalTok{ circuit.observables]}
            \ControlFlowTok{return}\NormalTok{ QuantumResult(expectation\_values}\OperatorTok{=}\NormalTok{evs, ...)}
        \ControlFlowTok{else}\NormalTok{:}
            \CommentTok{\# Sampler path → shot counts / bitstrings}
\NormalTok{            counts }\OperatorTok{=}\NormalTok{ sample\_bitstrings(circuit, options.shots)}
            \ControlFlowTok{return}\NormalTok{ QuantumResult(shots}\OperatorTok{=}\NormalTok{ShotResult(..., counts}\OperatorTok{=}\NormalTok{counts), ...)}
\end{Highlighting}
\end{Shaded}

\textbf{Rule}: SDK imports go inside the method that uses them --- never
at module level.
\end{frame}

\begin{frame}[fragile]{AlgorithmAdapter (QForte example)}
\phantomsection\label{algorithmadapter-qforte-example}
\begin{Shaded}
\begin{Highlighting}[]
\KeywordTok{class}\NormalTok{ QForteAlgorithmAdapter:}
    \KeywordTok{def}\NormalTok{ validate\_problem(}\VariableTok{self}\NormalTok{, circuit: CircuitSpec) }\OperatorTok{{-}\textgreater{}} \BuiltInTok{list}\NormalTok{[}\BuiltInTok{str}\NormalTok{]:}
        \ControlFlowTok{if}\NormalTok{ circuit.}\BuiltInTok{format} \OperatorTok{!=}\NormalTok{ CircuitFormat.MOLECULE\_JSON:}
            \ControlFlowTok{return}\NormalTok{ [}\SpecialStringTok{f"Expected MOLECULE\_JSON, got }\SpecialCharTok{\{}\NormalTok{circuit}\SpecialCharTok{.}\BuiltInTok{format}\SpecialCharTok{!r\}}\SpecialStringTok{"}\NormalTok{]}
        \ControlFlowTok{return}\NormalTok{ []}

    \KeywordTok{def}\NormalTok{ run\_algorithm(}\VariableTok{self}\NormalTok{, circuit: CircuitSpec,}
\NormalTok{                      options: ExecutionOptions}
\NormalTok{                      ) }\OperatorTok{{-}\textgreater{}} \BuiltInTok{tuple}\NormalTok{[QuantumResult, VQAConfig]:}
        \ImportTok{import}\NormalTok{ qforte }\ImportTok{as}\NormalTok{ qf                   }\CommentTok{\# ← deferred SDK import}
\NormalTok{        mol }\OperatorTok{=}\NormalTok{ qf.system\_factory(...)}
\NormalTok{        alg }\OperatorTok{=}\NormalTok{ qf.ADAPTVQE(mol, ...)}
\NormalTok{        alg.run(pool\_type}\OperatorTok{=}\StringTok{"SD"}\NormalTok{, optimizer}\OperatorTok{=}\StringTok{"BFGS"}\NormalTok{, ...)}
\NormalTok{        energy }\OperatorTok{=} \BuiltInTok{float}\NormalTok{(alg.get\_gs\_energy())}
\NormalTok{        result }\OperatorTok{=}\NormalTok{ QuantumResult(expectation\_values}\OperatorTok{=}\NormalTok{[}
\NormalTok{            ExpectationResult(observable\_index}\OperatorTok{=}\DecValTok{0}\NormalTok{, value}\OperatorTok{=}\NormalTok{energy, std\_error}\OperatorTok{=}\FloatTok{0.0}\NormalTok{)}
\NormalTok{        ], status}\OperatorTok{=}\NormalTok{JobStatus.SUCCEEDED)}
        \ControlFlowTok{return}\NormalTok{ result, VQAConfig(problem\_type}\OperatorTok{=}\StringTok{"chemistry"}\NormalTok{,}
\NormalTok{                                 final\_eigenvalue}\OperatorTok{=}\NormalTok{energy, ...)}
\end{Highlighting}
\end{Shaded}
\end{frame}

\begin{frame}[fragile]{Energy Evaluator Hook}
\phantomsection\label{energy-evaluator-hook}
Advanced: QForte drives the ansatz; a qpubench backend evaluates
\(\langle H \rangle\)

\begin{verbatim}
QForte optimizer  -->  energy_feval(params)
                              |
                    EnergyEvaluatorHook.evaluate()
                              |
                    qpubench BackendAdapter.run()
                              |
                    QuantumResult  -->  <H> (float)
\end{verbatim}

\begin{itemize}
\tightlist
\item
  Enables ADAPT-VQE on \textbf{Aer, Qrack GPU, IBM hardware}, or any
  adapter
\item
  Gradient computation automatically routed through the hook
\item
  Falls back to QForte's internal C++ simulator on error
\item
  Available in \texttt{integrations/qforte/energy\_hook.py}
\end{itemize}
\end{frame}

\begin{frame}[fragile]{Error Mitigation Providers}
\phantomsection\label{error-mitigation-providers}
\texttt{error\_mitigation.py} --- one schema module, seven provider
integrations:

\begin{longtable}[]{@{}
  >{\raggedright\arraybackslash}p{(\columnwidth - 4\tabcolsep) * \real{0.3333}}
  >{\raggedright\arraybackslash}p{(\columnwidth - 4\tabcolsep) * \real{0.3333}}
  >{\raggedright\arraybackslash}p{(\columnwidth - 4\tabcolsep) * \real{0.3333}}@{}}
\toprule\noalign{}
\begin{minipage}[b]{\linewidth}\raggedright
Provider
\end{minipage} & \begin{minipage}[b]{\linewidth}\raggedright
Schema types
\end{minipage} & \begin{minipage}[b]{\linewidth}\raggedright
Technique
\end{minipage} \\
\midrule\noalign{}
\endhead
\textbf{Q-CTRL Fire Opal} & \texttt{FireOpalConfig},
\texttt{FireOpalResult} & Closed-box noise suppression;
\texttt{fo.execute()} API \\
\textbf{Mitiq} & \texttt{MitiqConfig} + 5 configs, \texttt{MitiqResult}
& ZNE · PEC · CDR · REM · DDD \\
\textbf{Haiqu Rivet} & \texttt{HaiquRivetConfig},
\texttt{HaiquTranspilationResult} & Hardware-aware transpilation
caching \\
\textbf{ParityQC} & \texttt{ParityQCConfig}, \texttt{ParityQCResult} &
Parity encoding for QUBO / HCBO / Ising \\
\textbf{QMatter} & \texttt{QMatterConfig},
\texttt{QMatterCompressionResult} & Quantum problem compression \\
\textbf{Quantum Motion} & \texttt{QuantumMotionDeviceSpec} & Silicon
CMOS spin-qubit characterisation \\
\bottomrule\noalign{}
\end{longtable}

\vspace{0.3em}

\texttt{ErrorMitigationStrategy} enum extended: \texttt{FIRE\_OPAL} ·
\texttt{MITIQ\_ZNE/PEC/CDR/REM/DDD} · \texttt{HAIQU} ·
\texttt{PARITY\_QC} · \texttt{QMATTER}
\end{frame}

\begin{frame}[fragile]{IQM Hardware --- Star Topology}
\phantomsection\label{iqm-hardware-star-topology}
\begin{columns}[T]
\begin{column}{0.55\textwidth}
\textbf{IQM Resonance cloud} (\texttt{iqm-client} + \texttt{qiskit-iqm})

\begin{Shaded}
\begin{Highlighting}[]
\CommentTok{\# Star architecture: all qubits via COMPR1 resonator}
\NormalTok{backend }\OperatorTok{=}\NormalTok{ BackendSpec.iqm\_resonance(}
    \StringTok{"garnet"}\NormalTok{,       }\CommentTok{\# 20Q IQM Garnet}
\NormalTok{    api\_token\_ref}\OperatorTok{=}\StringTok{"IQM\_TOKEN"}\NormalTok{,}
\NormalTok{)}
\CommentTok{\# BackendSpec.iqm\_resonance("deneb")   \# 6Q}
\CommentTok{\# BackendSpec.iqm\_resonance("sirius")  \# 24Q Star}

\NormalTok{adapter }\OperatorTok{=}\NormalTok{ IQMAdapter(}\StringTok{"garnet"}\NormalTok{, token}\OperatorTok{=}\NormalTok{os.environ[}\StringTok{"IQM\_TOKEN"}\NormalTok{])}
\end{Highlighting}
\end{Shaded}

Calibration via
\texttt{IQMClient.get\_dynamic\_quantum\_architecture()}: T1
\textasciitilde28 µs · T2* \textasciitilde17 µs · PRX
\textasciitilde99.79\% · CZ \textasciitilde98.42\%
\end{column}

\begin{column}{0.45\textwidth}
\textbf{Native gate set}

\begin{longtable}[]{@{}
  >{\raggedright\arraybackslash}p{(\columnwidth - 4\tabcolsep) * \real{0.2500}}
  >{\raggedright\arraybackslash}p{(\columnwidth - 4\tabcolsep) * \real{0.4583}}
  >{\raggedright\arraybackslash}p{(\columnwidth - 4\tabcolsep) * \real{0.2917}}@{}}
\toprule\noalign{}
\begin{minipage}[b]{\linewidth}\raggedright
Gate
\end{minipage} & \begin{minipage}[b]{\linewidth}\raggedright
Parameters
\end{minipage} & \begin{minipage}[b]{\linewidth}\raggedright
Notes
\end{minipage} \\
\midrule\noalign{}
\endhead
\texttt{PRX} & \(\theta\), \(\varphi\) & Angles in \textbf{fractions of
turns} --- \(\theta\)=0.5 \(\rightarrow\) \(\pi\) rotation \\
\texttt{CZ} & --- & Controlled-Z, symmetric \\
\texttt{MOVE} & --- & Qubit \(\leftrightarrow\) resonator swap (Star
only) \\
\bottomrule\noalign{}
\end{longtable}

Hub-and-spoke topology gives effective all-to-all connectivity without
SWAP routing.
\end{column}
\end{columns}
\end{frame}

\begin{frame}[fragile]{IBM Quantum Runtime V2}
\phantomsection\label{ibm-quantum-runtime-v2}
\begin{columns}[T]
\begin{column}{0.5\textwidth}
\textbf{EstimatorV2 PUB}
\texttt{(circuit,\ observables,\ params?,\ precision?)}

\begin{Shaded}
\begin{Highlighting}[]
\CommentTok{\# precision replaces fixed shot count}
\NormalTok{estimator.run([(qc, obs, params, }\FloatTok{1e{-}2}\NormalTok{)])}
\NormalTok{result[i].data.evs   }\CommentTok{\# expectation value}
\NormalTok{result[i].data.stds  }\CommentTok{\# statistical error}
\end{Highlighting}
\end{Shaded}

\textbf{SamplerV2 PUB} \texttt{(circuit,\ params?,\ shots?)}

\begin{Shaded}
\begin{Highlighting}[]
\NormalTok{sampler.run([(qc, params, }\DecValTok{4096}\NormalTok{)])}
\NormalTok{bit\_array }\OperatorTok{=}\NormalTok{ result[}\DecValTok{0}\NormalTok{].data.meas}
\NormalTok{counts }\OperatorTok{=}\NormalTok{ bit\_array.get\_counts()}
\end{Highlighting}
\end{Shaded}

\texttt{BitArray.shape} = \texttt{(num\_shots,)} or
\texttt{(num\_params,\ num\_shots)}
\end{column}

\begin{column}{0.5\textwidth}
\textbf{Execution modes} (\texttt{IBMExecutionMode})

\begin{longtable}[]{@{}ll@{}}
\toprule\noalign{}
Mode & When to use \\
\midrule\noalign{}
\endhead
\texttt{SESSION} & VQE / adaptive --- exclusive QPU hold \\
\texttt{BATCH} & Sweeps --- parallel independent jobs \\
\texttt{SINGLE} & One-shot run (default) \\
\bottomrule\noalign{}
\end{longtable}

\begin{Shaded}
\begin{Highlighting}[]
\NormalTok{adapter }\OperatorTok{=}\NormalTok{ IBMAdapter(}
    \StringTok{"ibm\_torino"}\NormalTok{,}
\NormalTok{    execution\_mode}\OperatorTok{=}\NormalTok{IBMExecutionMode.SESSION,}
\NormalTok{)}
\end{Highlighting}
\end{Shaded}

\textbf{\texttt{IBMRuntimeRecord}} on \texttt{QuantumResult}: job\_id ·
session\_id · resilience\_level · \texttt{ExecutionSpan} timestamps

\vspace{0.3em}

\textbf{CUDA-Q}: \texttt{cudaq.observe()} \(\approx\) EstimatorV2;
\texttt{cudaq.sample()} \(\approx\) SamplerV2 --- same
\texttt{BackendAdapter} protocol.
\end{column}
\end{columns}
\end{frame}

\begin{frame}[fragile]{Quantum Advantage Tracker}
\phantomsection\label{quantum-advantage-tracker}
\texttt{BenchmarkRecord.advantage} holds a
\texttt{QuantumAdvantageRecord}:

\begin{Shaded}
\begin{Highlighting}[]
\NormalTok{record }\OperatorTok{=}\NormalTok{ runner.run(circuit, }\StringTok{"ibm\_torino"}\NormalTok{, options)}
\NormalTok{record }\OperatorTok{=}\NormalTok{ record.model\_copy(update}\OperatorTok{=}\NormalTok{\{}\StringTok{"advantage"}\NormalTok{: QuantumAdvantageRecord(}
\NormalTok{    experiment\_type}\OperatorTok{=}\NormalTok{AdvantageExperimentType.OBSERVABLE\_ESTIMATION,}
\NormalTok{    circuit\_name}\OperatorTok{=}\StringTok{"Loschmidt{-}echo{-}70q"}\NormalTok{,}
\NormalTok{    num\_qubits}\OperatorTok{=}\DecValTok{70}\NormalTok{,}
\NormalTok{    backend\_name}\OperatorTok{=}\StringTok{"ibm\_boston"}\NormalTok{,}
\NormalTok{    floquet\_layers}\OperatorTok{=}\DecValTok{20}\NormalTok{,}
\NormalTok{    observable\_value}\OperatorTok{=}\FloatTok{0.327}\NormalTok{,}
\NormalTok{    observable\_error\_bound}\OperatorTok{=}\FloatTok{0.012}\NormalTok{,}
\NormalTok{    classical\_method}\OperatorTok{=}\NormalTok{ClassicalComparisonMethod.TENSOR\_NETWORK,}
\NormalTok{    coupling\_params}\OperatorTok{=}\NormalTok{\{}\StringTok{"b"}\NormalTok{: }\FloatTok{1.0}\NormalTok{, }\StringTok{"delta"}\NormalTok{: }\FloatTok{0.5}\NormalTok{\},}
\NormalTok{    submission\_url}\OperatorTok{=}\StringTok{"https://github.com/quantum{-}advantage{-}tracker/..."}\NormalTok{,}
\NormalTok{)\})}
\end{Highlighting}
\end{Shaded}

\begin{block}{Quantum Advantage Tracker}
Community registry co-initiated by IBM, Flatiron Institute, BlueQubit, Algorithmiq.
Three experiment categories: observable estimation · variational · classically verifiable.
\end{block}
\end{frame}

\section{Execution \& Storage}\label{execution-storage}

\begin{frame}[fragile]{BenchmarkRunner}
\phantomsection\label{benchmarkrunner}
\begin{Shaded}
\begin{Highlighting}[]
\NormalTok{runner }\OperatorTok{=}\NormalTok{ BenchmarkRunner(store}\OperatorTok{=}\NormalTok{NDJSONStore(Path(}\StringTok{"results.ndjson"}\NormalTok{)))}
\NormalTok{runner.register(AerAdapter(), name}\OperatorTok{=}\StringTok{"aer"}\NormalTok{)}
\NormalTok{runner.register(QForteAlgorithmAdapter(), name}\OperatorTok{=}\StringTok{"qforte"}\NormalTok{)}

\CommentTok{\# Single run}
\NormalTok{record }\OperatorTok{=}\NormalTok{ runner.run(circuit, }\StringTok{"aer"}\NormalTok{, ExecutionOptions(shots}\OperatorTok{=}\DecValTok{4096}\NormalTok{))}

\CommentTok{\# Cartesian product sweep: circuits × backends × options}
\NormalTok{records }\OperatorTok{=}\NormalTok{ runner.sweep(}
\NormalTok{    circuits}\OperatorTok{=}\NormalTok{[circuit\_h2, circuit\_lih],}
\NormalTok{    backend\_names}\OperatorTok{=}\NormalTok{[}\StringTok{"aer"}\NormalTok{, }\StringTok{"qforte"}\NormalTok{],}
\NormalTok{    options\_list}\OperatorTok{=}\NormalTok{[ExecutionOptions(shots}\OperatorTok{=}\DecValTok{1024}\NormalTok{),}
\NormalTok{                  ExecutionOptions(shots}\OperatorTok{=}\DecValTok{8192}\NormalTok{)],}
\NormalTok{    run\_id}\OperatorTok{=}\StringTok{"vqe{-}sweep{-}001"}\NormalTok{,}
\NormalTok{)}

\CommentTok{\# Post{-}execution hooks}
\NormalTok{runner.add\_hook(}\KeywordTok{lambda}\NormalTok{ rec: log\_energy(rec.vqa.final\_eigenvalue))}
\end{Highlighting}
\end{Shaded}
\end{frame}

\begin{frame}[fragile]{Data Persistence}
\phantomsection\label{data-persistence}
\textbf{NDJSONStore} --- zero dependencies, append-only, grep-able

\begin{Shaded}
\begin{Highlighting}[]
\NormalTok{store }\OperatorTok{=}\NormalTok{ NDJSONStore(Path(}\StringTok{"results/run1.ndjson"}\NormalTok{))}
\NormalTok{store.save(record)}

\NormalTok{record  }\OperatorTok{=}\NormalTok{ store.load(experiment\_id)           }\CommentTok{\# by UUID}
\NormalTok{records }\OperatorTok{=}\NormalTok{ store.query(                        }\CommentTok{\# dot{-}path filters}
\NormalTok{    backend\_\_name}\OperatorTok{=}\StringTok{"aer\_statevector"}\NormalTok{,}
\NormalTok{    result\_\_status}\OperatorTok{=}\StringTok{"succeeded"}\NormalTok{,}
\NormalTok{)}
\end{Highlighting}
\end{Shaded}

\textbf{ParquetStore} --- columnar analytics
(\texttt{pip\ install\ ".{[}storage{]}"})

\begin{Shaded}
\begin{Highlighting}[]
\NormalTok{store }\OperatorTok{=}\NormalTok{ ParquetStore(Path(}\StringTok{"results/run1.parquet"}\NormalTok{))}
\NormalTok{df }\OperatorTok{=}\NormalTok{ store.to\_dataframe()                     }\CommentTok{\# → pandas DataFrame}
\end{Highlighting}
\end{Shaded}

Every record is one JSON line --- stream-able, pandas-ready,
archive-safe.
\end{frame}

\section{Integrations}\label{integrations}

\begin{frame}[fragile]{Integration Ecosystem (1/3)}
\phantomsection\label{integration-ecosystem-13}
\begin{longtable}[]{@{}
  >{\raggedright\arraybackslash}p{(\columnwidth - 4\tabcolsep) * \real{0.3611}}
  >{\raggedright\arraybackslash}p{(\columnwidth - 4\tabcolsep) * \real{0.3889}}
  >{\raggedright\arraybackslash}p{(\columnwidth - 4\tabcolsep) * \real{0.2500}}@{}}
\toprule\noalign{}
\begin{minipage}[b]{\linewidth}\raggedright
Integration
\end{minipage} & \begin{minipage}[b]{\linewidth}\raggedright
Schema module
\end{minipage} & \begin{minipage}[b]{\linewidth}\raggedright
Protocol
\end{minipage} \\
\midrule\noalign{}
\endhead
Cebule SDK & \texttt{cebule} & \texttt{BackendAdapter} \\
Xenakis GA circuit search & \texttt{xenakis} &
\texttt{AlgorithmAdapter} \\
ExcitationSolve (Fourier VQE) & \texttt{excitation\_solve} &
\texttt{AlgorithmAdapter} \\
GSOpt (ground-state benchmark) & \texttt{gsopt} &
\texttt{AlgorithmAdapter} \\
Photochipsim / FBQC & \texttt{photonic} & \texttt{BackendAdapter} \\
QDK / QuNorth chemistry & \texttt{qdk\_chemistry} &
\texttt{AlgorithmAdapter} \\
\bottomrule\noalign{}
\end{longtable}
\end{frame}

\begin{frame}[fragile]{Integration Ecosystem (2/3)}
\phantomsection\label{integration-ecosystem-23}
\begin{longtable}[]{@{}lll@{}}
\toprule\noalign{}
Integration & Schema module & Protocol \\
\midrule\noalign{}
\endhead
DTU-GBS / photonic\_QC & \texttt{gbs} & \texttt{BackendAdapter} \\
QSE / KQD & \texttt{qse} & \texttt{AlgorithmAdapter} \\
QESEM (Qedma) & \texttt{qesem} & \texttt{BackendAdapter} \\
QCSchema / PennyLane & \texttt{qcschema} & cross-format interop \\
Bloqade / Aquila & \texttt{neutral\_atom} & \texttt{BackendAdapter} \\
SlowQuant UCC/VQE & \texttt{slowquant} & \texttt{AlgorithmAdapter} \\
\bottomrule\noalign{}
\end{longtable}
\end{frame}

\begin{frame}[fragile]{Integration Ecosystem (3/3) --- Error Mitigation
\& Hardware}
\phantomsection\label{integration-ecosystem-33-error-mitigation-hardware}
\begin{longtable}[]{@{}
  >{\raggedright\arraybackslash}p{(\columnwidth - 4\tabcolsep) * \real{0.3611}}
  >{\raggedright\arraybackslash}p{(\columnwidth - 4\tabcolsep) * \real{0.3889}}
  >{\raggedright\arraybackslash}p{(\columnwidth - 4\tabcolsep) * \real{0.2500}}@{}}
\toprule\noalign{}
\begin{minipage}[b]{\linewidth}\raggedright
Integration
\end{minipage} & \begin{minipage}[b]{\linewidth}\raggedright
Schema module
\end{minipage} & \begin{minipage}[b]{\linewidth}\raggedright
Protocol
\end{minipage} \\
\midrule\noalign{}
\endhead
Q-CTRL Fire Opal & \texttt{error\_mitigation} &
\texttt{ErrorMitigationAdapter} \\
Mitiq (ZNE · PEC · CDR · REM · DDD) & \texttt{error\_mitigation} &
\texttt{ErrorMitigationAdapter} \\
Haiqu Rivet transpiler & \texttt{error\_mitigation} &
\texttt{ErrorMitigationAdapter} \\
ParityQC parity encoding & \texttt{error\_mitigation} &
\texttt{ErrorMitigationAdapter} \\
QMatter problem compression & \texttt{error\_mitigation} &
\texttt{ErrorMitigationAdapter} \\
IQM Resonance (garnet · deneb · sirius) & --- &
\texttt{BackendAdapter} \\
Quantum Motion CMOS spin-qubit & \texttt{error\_mitigation} &
\texttt{BackendAdapter} \\
\bottomrule\noalign{}
\end{longtable}
\end{frame}

\begin{frame}[fragile]{Cross-Framework Compatibility}
\phantomsection\label{cross-framework-compatibility}
\textbf{Pauli encoding} --- explicit converters prevent silent
convention bugs:

\begin{Shaded}
\begin{Highlighting}[]
\CommentTok{\# Qrack Q\# convention: I=0, X=1, Z=2, Y=3  (non{-}sequential!)}
\NormalTok{PauliLabel.Z.to\_qrack\_int()           }\CommentTok{\# → 2  (not 3)}
\NormalTok{PauliLabel.Y.to\_qrack\_int()           }\CommentTok{\# → 3  (not 2)}

\CommentTok{\# Qiskit C API QkBitTerm bit{-}packed encoding}
\NormalTok{PauliLabel.X.to\_qiskit\_c\_bit\_term()   }\CommentTok{\# → 0b0010}

\CommentTok{\# MBQC byproduct register: bit 0 = Z, bit 1 = X  (reversed vs gate{-}based!)}
\end{Highlighting}
\end{Shaded}

\textbf{Complex numbers}: stored as \texttt{\{"re":\ 1.0,\ "im":\ 0.5\}}
--- valid JSON without Python eval

\textbf{QCSchema / QCElemental / PennyLane} interoperability via the
\texttt{qcschema} module
\end{frame}

\section{Code Quality}\label{code-quality}

\begin{frame}[fragile]{Python Quality Toolchain}
\phantomsection\label{python-quality-toolchain}
\begin{Shaded}
\begin{Highlighting}[]
\ExtensionTok{ruff}\NormalTok{ check src/ tests/    }\CommentTok{\# lint (line{-}length 100, target py311)}
\ExtensionTok{ruff}\NormalTok{ format src/ tests/   }\CommentTok{\# format}
\ExtensionTok{mypy}\NormalTok{ src/                 }\CommentTok{\# static types (strict = true)}
\ExtensionTok{pytest}\NormalTok{ tests/ }\AttributeTok{{-}q}          \CommentTok{\# 156 tests — no quantum SDK required}
\end{Highlighting}
\end{Shaded}

\begin{longtable}[]{@{}
  >{\raggedright\arraybackslash}p{(\columnwidth - 2\tabcolsep) * \real{0.5000}}
  >{\raggedright\arraybackslash}p{(\columnwidth - 2\tabcolsep) * \real{0.5000}}@{}}
\toprule\noalign{}
\begin{minipage}[b]{\linewidth}\raggedright
Tool
\end{minipage} & \begin{minipage}[b]{\linewidth}\raggedright
Role
\end{minipage} \\
\midrule\noalign{}
\endhead
\textbf{Pydantic v2} & Runtime-validated schemas --- field types
enforced at parse time \\
\textbf{ruff} & Linting + auto-formatting, replaces
flake8/isort/black \\
\textbf{mypy strict} & Full type coverage --- \texttt{Any} only where
unavoidable \\
\textbf{pytest} & Schema-only test suite runs with
\texttt{pip\ install\ .} alone \\
\textbf{ponytail} & Agent code-quality enforcement via decision
ladder \\
\bottomrule\noalign{}
\end{longtable}
\end{frame}

\section{Summary}\label{summary}

\begin{frame}[fragile]{Getting Started}
\phantomsection\label{getting-started}
\begin{Shaded}
\begin{Highlighting}[]
\ImportTok{from}\NormalTok{ qpubench }\ImportTok{import}\NormalTok{ (}
\NormalTok{    BenchmarkRunner, NDJSONStore, StubGateAdapter,}
\NormalTok{    CircuitSpec, ExecutionOptions,}
\NormalTok{    SparsePauliObservable, PauliTerm, PauliLabel, ComplexNumber,}
\NormalTok{)}
\ImportTok{import}\NormalTok{ pathlib}

\NormalTok{circuit }\OperatorTok{=}\NormalTok{ CircuitSpec(}
\NormalTok{    num\_qubits}\OperatorTok{=}\DecValTok{2}\NormalTok{,}
\NormalTok{    serialized}\OperatorTok{=}\StringTok{\textquotesingle{}OPENQASM 2.0;}\CharTok{\textbackslash{}n}\StringTok{include "qelib1.inc";}\CharTok{\textbackslash{}n}\StringTok{qreg q[2];}\CharTok{\textbackslash{}n}\StringTok{h q[0];}\CharTok{\textbackslash{}n}\StringTok{cx q[0],q[1];\textquotesingle{}}\NormalTok{,}
\NormalTok{    observables}\OperatorTok{=}\NormalTok{[SparsePauliObservable(num\_qubits}\OperatorTok{=}\DecValTok{2}\NormalTok{, terms}\OperatorTok{=}\NormalTok{[}
\NormalTok{        PauliTerm(qubit\_indices}\OperatorTok{=}\NormalTok{(}\DecValTok{0}\NormalTok{, }\DecValTok{1}\NormalTok{),}
\NormalTok{                  pauli\_ops}\OperatorTok{=}\NormalTok{(PauliLabel.Z, PauliLabel.Z),}
\NormalTok{                  coefficient}\OperatorTok{=}\NormalTok{ComplexNumber(re}\OperatorTok{=}\FloatTok{1.0}\NormalTok{))}
\NormalTok{    ])],}
\NormalTok{)}
\NormalTok{runner }\OperatorTok{=}\NormalTok{ BenchmarkRunner(store}\OperatorTok{=}\NormalTok{NDJSONStore(pathlib.Path(}\StringTok{"results.ndjson"}\NormalTok{)))}
\NormalTok{runner.register(StubGateAdapter(seed}\OperatorTok{=}\DecValTok{42}\NormalTok{), name}\OperatorTok{=}\StringTok{"stub"}\NormalTok{)}
\NormalTok{record }\OperatorTok{=}\NormalTok{ runner.run(circuit, }\StringTok{"stub"}\NormalTok{, ExecutionOptions(shots}\OperatorTok{=}\DecValTok{4096}\NormalTok{))}
\BuiltInTok{print}\NormalTok{(}\SpecialStringTok{f"\textless{}ZZ\textgreater{} = }\SpecialCharTok{\{}\NormalTok{record}\SpecialCharTok{.}\NormalTok{result}\SpecialCharTok{.}\NormalTok{expectation\_values[}\DecValTok{0}\NormalTok{]}\SpecialCharTok{.}\NormalTok{value}\SpecialCharTok{:.4f\}}\SpecialStringTok{"}\NormalTok{)}
\end{Highlighting}
\end{Shaded}
\end{frame}

\begin{frame}[fragile]{Roadmap \& Links}
\phantomsection\label{roadmap-links}
\textbf{Repository}: \texttt{github.com/mqsdk/qpubench}

\textbf{Schema v1.12.0} --- 21 modules · 9 QPU modalities · zero quantum
SDK deps in core

\vspace{0.5em}

\textbf{Contributing} --- writing an adapter takes \textasciitilde30
lines:

\begin{Shaded}
\begin{Highlighting}[]
\FunctionTok{cp}\NormalTok{ integrations/template/backend\_adapter\_template.py  my\_adapter.py}
\CommentTok{\# fill the three TODOs: spec, validate(), run()}
\end{Highlighting}
\end{Shaded}

See \texttt{INTEGRATION\_GUIDE.md} and \texttt{AGENTS.md} for the full
development workflow.

\vspace{0.5em}

\textbf{Contact}: \texttt{contact@mqs.dk}
\end{frame}

\end{document}
